Surface operations on the Moon will depend on short-range links between assets
that are close together but separately powered: a rover and its lander, a
sensor mast and a relay node, an instrument package and the vehicle that
deployed it. These \emph{proximity} links carry housekeeping telemetry,
commands, and instrument data over metres rather than kilometres, and they are
not the same design problem as the high-rate optical trunk links that dominate
the space optical communication literature.

Optical signalling is attractive at this scale for reasons specific to the
environment. It requires no spectrum allocation and generates no
radio-frequency emissions that might couple into sensitive instrumentation. It
is inherently directional, which limits mutual interference between co-located
assets. Emitter and photodiode front ends are light and draw little power,
which matters where every gram and every watt is contested. Permanently
shadowed regions near the lunar poles, of particular interest for volatiles
prospecting, additionally offer very low ambient infrared background, an
operating condition that favours optical reception and is difficult to
reproduce terrestrially.

The constraint that shapes such a link is not aperture or transmit power but
the receiver. A surface asset's communication processor is small,
radiation-tolerant, and shared with other tasks, so decoding must be cheap
enough to run inside a sampling loop. In practice this means threshold
decoding: each detected pulse is assigned to a symbol by testing it against
fixed acceptance windows. That approach has a failure mode which is rarely
quantified. It has no representation for a pulse that was never detected at
all, so when the receiver misses a symbol the byte containing it is lost, and
because protocols of this kind spend several pulses per byte, a small
per-symbol error rate compounds into a large message error rate. Improving the
per-symbol classifier cannot help, because the missing observation was never
presented to it.

Message loss is more costly in this setting than the raw numbers suggest.
Proximity links operate on constrained duty cycles and within scheduled
communication windows, so a lost message is not simply retransmitted at
negligible cost: retransmission consumes energy, occupies window time, and on a
store-and-forward relay path can delay delivery by an entire orbit. Recovering
a message at the receiver, using information already present in the captured
signal, is therefore worth more than the same improvement would be on a
terrestrial link with cheap retries.

This paper presents LORAD (Lunar Optical Relay And Decoding), an experimental
platform built to measure that effect and to address it. The link uses baseband
on-off keying in which symbol identity is carried by pulse \emph{width} rather
than amplitude. That choice is not incidental. Measurements presented in
Section~\ref{results} show the receiver's transimpedance stage to be saturated,
pinning received peak amplitude to within 0.17\% across every pulse recorded,
so an amplitude-discriminating scheme could not function on this hardware,
whereas pulse-width classes remain separated by 12--40 pooled standard
deviations even at rates where decoding fails outright. Time-domain
encoding is also the more defensible choice for a dusty, thermally cycling
surface environment, where optical throughput and detector responsivity drift
but symbol timing does not.

The receiver streams per-pulse observations to a decoder that treats decoding
as sequence estimation rather than classification. A hidden Markov model
encodes the protocol's byte framing directly in its state, so that transitions
producing an inadmissible byte are pruned during search rather than filtered
afterwards. Skip transitions represent undetected symbols, and inter-pulse gap
timing supplies the evidence that makes them identifiable: a symbol that goes
undetected lengthens the following gap by its own width plus one inter-symbol
gap, indicating both \emph{that} a symbol was lost and \emph{which} one. No
change to the transmitter, the receiver electronics, or the modulation is
required: the improvement is obtained entirely from information the existing
receiver already captures and discards.

The contributions of this work are:

\begin{itemize}
\item A characterization of a pulse-width OOK optical proximity link across
five data rates from 28.6 to 171.4~bit/s, establishing that its dominant
failure mode is undetected pulses rather than misclassified ones, and that
every such omission nonetheless leaves a recoverable signature in the
inter-pulse timing. This yields a reference level for what a decoder operating
on the detected pulse stream can be expected to recover.

\item A maximum-likelihood sequence decoder combining a learned emission model,
gap-timing evidence, and the protocol's framing and printable-ASCII
constraints, which raises transmission success from 72.2\% to 100\% at 28.6~bit/s and
from 36.1\% to 69.4\% at 114.3~bit/s on held-out data.

\item A live closed-loop demonstration in which a memory-constrained receiver
streams observations to the decoder and receives corrected messages back,
raising transmission success from 5.9\% to 44.9\% at 114.3~bit/s over 118
transmissions, an architecture directly analogous to a surface asset
offloading decoding to a better-resourced relay or lander.

\item An account of the measurement pitfalls encountered on a
resource-constrained receiver: an instrumentation artifact that inflated
apparent symbol loss by as much as a factor of two and was mistaken for channel
degradation until isolated, and a payload-length dependence that accounted for
most of an apparent discrepancy between two experiments. A by-product of the
first is a method for recovering the receiver's detection-loop period, and with
it a fixed edge-detection bias, from the periodicity of its own pulse-width
measurements.
\end{itemize}

The remainder of this paper is organized as follows.
Section~\ref{related_works} reviews prior work in space optical communication,
proximity links, receiver diversity, and sequence detection.
Section~\ref{system_overview} describes the link protocol, receiver hardware,
and decoder model. Section~\ref{experiments} details the experimental method
and datasets. Section~\ref{results} presents the measurements, and
Section~\ref{conclusion} concludes.
